\documentclass[12pt,a4paper]{article}

% ─── Packages ────────────────────────────────────────────────
\usepackage{amsmath,amssymb}
\usepackage{booktabs}
\usepackage[margin=1in]{geometry}
\usepackage{setspace}
\usepackage{hyperref}
\usepackage{xcolor}
\usepackage{float}
\usepackage{fontspec}
\usepackage{xeCJK}
\setCJKmainfont{PingFang TC}
\usepackage{enumitem}
\usepackage{longtable}
\usepackage{array}

\hypersetup{colorlinks=true, linkcolor=blue, citecolor=blue, urlcolor=blue}
\onehalfspacing

% Colors for diff marking
\newcommand{\added}[1]{\textcolor{blue}{#1}}
\newcommand{\removed}[1]{\textcolor{red}{\sout{#1}}}
\newcommand{\changed}[1]{\textcolor{orange}{#1}}
\usepackage[normalem]{ulem}  % for \sout

\title{Revision Diff Report: v1 $\rightarrow$ v2\\
\large ``Can Anything Beat VIX?''\\[6pt]
\normalsize Addressing TOP 3 HIGH-Severity Review Issues + Citation Fixes}

\author{VolPred Research System}
\date{March 30, 2026}

\begin{document}

\maketitle

%=================================================================
\section{Overview}
%=================================================================

\begin{tabular}{@{}lcc@{}}
\toprule
\textbf{Metric} & \textbf{v1} & \textbf{v2} \\
\midrule
Pages & 31 & 34 \\
References & 24 & 33 \\
$R^2$ metrics reported & 1 (VIX-relative only) & 2 (Campbell-Thompson + VIX-relative) \\
Multiple-testing method & $|t|>3.0$ rule of thumb & $|t|>3.0$ + Holm-Bonferroni FWER \\
Table 2 columns & 7 & 10 \\
Issues addressed & --- & 3 HIGH + 3 citation fixes \\
\bottomrule
\end{tabular}

\bigskip

\noindent\textbf{Page count change}: +3 pages (31 $\rightarrow$ 34). The increase comes from:
\begin{itemize}
    \item Expanded literature review in Section 2 (+1 page): Poon \& Granger, Andersen et al., Hansen \& Lunde, Patton \& Sheppard, Patton (2011)
    \item Dual $R^2$ definition in Section 4 (+0.5 page): Campbell-Thompson formula + explanation paragraph
    \item Expanded multiple-testing section in Section 5 (+0.5 page): Holm-Bonferroni procedure description
    \item Wider Table 2 with 3 new columns (+0.5 page): $R^2_{\text{OOS,CT}}$, Raw $p$, HB $p$
    \item New discussion paragraphs after Table 2 (+0.5 page): dual $R^2$ interpretation + Holm-Bonferroni results
    \item 9 new bibliography entries (+0.5 page)
\end{itemize}


%=================================================================
\section{Issue D2.1: Missing Foundational Literature [HIGH]}
%=================================================================

\subsection*{Problem}
Eight foundational papers were not cited: Poon \& Granger (2003), Hansen \& Lunde (2005), Andersen et al.\ (2003), Patton \& Sheppard (2015), Campbell \& Thompson (2008), Newey \& West (1987), Patton (2011), Holm (1979).

\subsection*{Changes Made}

\begin{enumerate}

\item \textbf{Section 2.2 (Empirical Evidence)} --- Added opening paragraph:

\added{``The volatility forecasting literature is vast; Poon and Granger (2003) provide the definitive survey, cataloguing over 90 studies and concluding that implied volatility generally dominates time-series models.''}

\added{Added new paragraph on Andersen et al.\ (2003) for realized volatility methodology and Patton (2011) for loss function choice.}

\item \textbf{Section 2.4 (Prior Horse Races)} --- Added opening paragraph:

\added{``The canonical volatility model horse race is Hansen and Lunde (2005), who compare 330 ARCH-type models for forecasting daily exchange rate and IBM stock volatility. Their central finding---that no model significantly outperforms the simple GARCH(1,1)---sets the precedent for our study\ldots''}

\added{Added Patton and Sheppard (2015) for good/bad volatility decomposition.}

\item \textbf{Section 5.1 (DM Test)} --- Added Newey-West citation:

\added{``\ldots with Newey and West (1987) heteroskedasticity and autocorrelation consistent (HAC) standard errors (bandwidth = 22, matching the forecast horizon).''}

\item \textbf{Bibliography} --- Added 9 new entries:

\begin{itemize}
    \item Andersen, Bollerslev, Diebold, and Labys (2003) -- \textit{Econometrica}
    \item Campbell and Thompson (2008) -- \textit{Review of Financial Studies}
    \item Engle and Rangel (2008) -- \textit{Review of Financial Studies}
    \item Hansen and Lunde (2005) -- \textit{Journal of Applied Econometrics}
    \item Holm (1979) -- \textit{Scandinavian Journal of Statistics}
    \item Lou, Polk, and Skouras (2019) -- \textit{Journal of Financial Economics}
    \item Newey and West (1987) -- \textit{Econometrica}
    \item Patton (2011) -- \textit{Journal of Econometrics}
    \item Patton and Sheppard (2015) -- \textit{Review of Economics and Statistics}
    \item Poon and Granger (2003) -- \textit{Journal of Economic Literature}
\end{itemize}

Total references: \removed{24} $\rightarrow$ \added{33} (+9 new, including 1 replacement).

\end{enumerate}


%=================================================================
\section{Issue D3.1: Non-Standard OOS $R^2$ [HIGH]}
%=================================================================

\subsection*{Problem}
OOS $R^2$ was defined relative to VIX forecasts (Eq.\ 5 in v1) rather than the Campbell-Thompson (2008) convention (relative to historical mean), making results incomparable to the existing literature.

\subsection*{Changes Made}

\begin{enumerate}

\item \textbf{Section 4.4 (Out-of-Sample $R^2$)} --- Completely rewritten to define \textbf{both} metrics:

\added{New Equation (5a): $R^2_{\text{OOS,CT}} = 1 - \frac{\sum (RV_t - \hat{RV}_t^{\text{model}})^2}{\sum (RV_t - \bar{RV}_{1:t-1})^2}$ \quad [Campbell-Thompson, relative to historical mean]}

\added{New Equation (5b): $R^2_{\text{OOS,VIX}} = 1 - \frac{\sum (RV_t - \hat{RV}_t^{\text{model}})^2}{\sum (RV_t - \hat{RV}_t^{\text{VIX}})^2}$ \quad [VIX-relative, our primary metric]}

\added{Added explanatory paragraph: ``We report both metrics to ensure comparability with the existing literature\ldots many signals that achieve positive $R^2_{\text{OOS,CT}}$ (i.e., beat the historical mean) still show zero $R^2_{\text{OOS,VIX}}$ (i.e., add nothing beyond VIX).''}

\item \textbf{Table 2} --- Added new column $R^2_{\text{OOS,CT}}$:

\added{All 8 testable signal families now report Campbell-Thompson $R^2_{\text{OOS,CT}}$ values (range: 0.312--0.398). The VIX-only benchmark achieves $R^2_{\text{OOS,CT}} = 0.365$.}

\item \textbf{Discussion after Table 2} --- Added new paragraph:

\added{``The dual $R^2$ metrics reveal a striking pattern: most signal families achieve positive Campbell-Thompson $R^2_{\text{OOS,CT}}$ values in the range of 0.31--0.40, indicating that they do forecast volatility better than the naive historical mean. However, these same signals show zero or negative VIX-relative $R^2_{\text{OOS,VIX}}$, confirming that their forecasting power is entirely subsumed by VIX.''}

\item \textbf{Abstract} --- Updated to mention both metrics:

\added{``\ldots whether measured by the standard Campbell and Thompson (2008) $R^2_{\text{OOS}}$ relative to the historical mean, by incremental $R^2$ relative to VIX-only forecasts, or by Sharpe ratio gains\ldots''}

\end{enumerate}


%=================================================================
\section{Issue D4.1: No Formal Multiple-Testing Correction [HIGH]}
%=================================================================

\subsection*{Problem}
Only $|t|>3.0$ rule of thumb from Harvey et al.\ (2016). No Bonferroni, Holm-Bonferroni, or BH FDR. Harvey et al.\ recommendation was oversimplified.

\subsection*{Changes Made}

\begin{enumerate}

\item \textbf{Section 5.2} --- Completely rewritten and retitled from ``Harvey (2016) Multiple-Testing Threshold'' to ``Multiple-Testing Corrections'':

\added{Clarified Harvey et al.\ (2016): ``Their analysis---applied to roughly 300 factors in the cross-section literature---yields an approximate threshold of $|t| > 3.0$ under conservative priors. We adopt $|t| > 3.0$ as an approximate conservative threshold motivated by their analysis, noting that for our 11 tests the simple Bonferroni correction yields a comparable threshold of $|t| > 3.23$.'' [Fixes citation check issue: harvey2016 oversimplified claim]}

\added{Added formal Holm-Bonferroni procedure: ``We implement the Holm-Bonferroni sequential procedure (Holm, 1979) across the 11 primary Diebold-Mariano tests. The procedure works as follows: (i) rank the 11 raw $p$-values from smallest to largest; (ii) for rank $k$, reject $H_0$ if $p_{(k)} \leq \alpha / (11 - k + 1)$; (iii) stop at the first non-rejection.''}

\item \textbf{Table 2} --- Added two new columns:

\added{Raw $p$: Two-sided $p$-values from DM tests (range: 0.147--0.881)}

\added{HB $p$: Holm-Bonferroni adjusted $p$-values (all 1.000, confirming null survives formal correction)}

\item \textbf{Introduction safeguard (iv)} --- Updated:

\removed{``We adopt the Harvey (2016) threshold of $|t| > 3.0$\ldots''}

\added{``We adopt the approximate conservative threshold of $|t| > 3.0$ motivated by Harvey et al.\ (2016) and implement formal Holm-Bonferroni sequential correction across all eleven Diebold-Mariano tests, controlling the family-wise error rate at $\alpha = 0.05$.''}

\item \textbf{Discussion after Table 2} --- Added new paragraph:

\added{``The Holm-Bonferroni correction reinforces the null. All eight testable DM tests produce raw $p$-values well above 0.05 (minimum raw $p = 0.147$), and the adjusted $p$-values are uniformly 1.000. The multiplicity correction is strictly redundant here\ldots but reporting it demonstrates that our conclusions are robust to formal multiple-testing control.''}

\item \textbf{Abstract and Conclusion} --- Updated to mention Holm-Bonferroni.

\end{enumerate}


%=================================================================
\section{Citation Fixes (from Citation Check Report)}
%=================================================================

\subsection{harvey2016: Oversimplified Claim [HIGH]}

\removed{``They recommend $|t| > 3.0$ as a minimum threshold for asset pricing tests.''}

\added{``Their analysis---applied to roughly 300 factors in the cross-section literature---yields an approximate threshold of $|t| > 3.0$ under conservative priors. We adopt $|t| > 3.0$ as an approximate conservative threshold motivated by their analysis.''}

\subsection{schwert1989: Misleading Positioning [HIGH]}

\removed{``economic indicators (Schwert, 1989)''}

\added{``macroeconomic determinants of time-varying volatility (Schwert, 1989; Engle and Rangel, 2008)'' --- Added Engle \& Rangel (2008) as the more directly relevant citation for economic indicator signals, while retaining Schwert for time-varying volatility documentation.}

\subsection{bollerslev2020: Content Mismatch [MEDIUM]}

\removed{``Overnight price changes reflect the processing of after-hours news\ldots (Bollerslev et al., 2020)'' --- Bollerslev et al.\ (2020) is about the cross-section of stock returns and semivariance, NOT overnight returns.}

\added{``Overnight price changes reflect the processing of after-hours news\ldots (Lou et al., 2019)'' --- Lou, Polk, and Skouras (2019) ``A Tug of War: Overnight Versus Intraday Expected Returns'' is the correct citation for overnight return dynamics.}

Note: The \texttt{bollerslev2020} bibitem is retained in the bibliography for potential use elsewhere but is no longer cited in the text for the overnight claim.


%=================================================================
\section{Additional Improvements}
%=================================================================

\begin{enumerate}
    \item \textbf{Table label fix}: Moved \texttt{\textbackslash caption} and \texttt{\textbackslash label} outside \texttt{threeparttable} for Tables 1 and 2 to resolve cross-reference warnings that existed in v1.

    \item \textbf{Table 2 expanded}: From 7 columns to 10 columns ($R^2_{\text{OOS,CT}}$, Raw $p$, HB $p$), with \texttt{resizebox} to fit within margins.

    \item \textbf{VIX-only benchmark row}: Added to Table 2 showing $R^2_{\text{OOS,CT}} = 0.365$ for the VIX benchmark itself, giving readers a reference point for signal comparison.
\end{enumerate}


%=================================================================
\section{Issues NOT Addressed in This Revision}
%=================================================================

The following review issues were not in scope for this revision but should be addressed in v3:

\begin{itemize}
    \item D3.2 [HIGH]: Uneven sample periods (811 to 8,325 days)
    \item D7.1 [HIGH]: Missing robustness checks (forecast horizon, QLIKE, international)
    \item D1.1 [MEDIUM]: ``Pre-registered'' $\rightarrow$ ``pre-specified''
    \item D2.2 [MEDIUM]: Post-Moreira-Muir debate
    \item D5.1 [MEDIUM]: Incomplete results for Families 5--7
    \item D5.2 [MEDIUM]: Inconsistent Table 3 numbering
    \item D6.1 [MEDIUM]: Proposition 1 approximation
    \item D9.1 [MEDIUM]: No MCS analysis
    \item D10.1 [MEDIUM]: Paper length
    \item All LOW issues
    \item \texttt{luo2019} [MEDIUM citation]: Weak support for contango/backwardation claim (not addressed --- requires finding a better primary citation)
\end{itemize}


\end{document}
